
\begin{lemma}\label{lemma:parametrization}
    Let $n = rs$, $K = K_{r,s} \cong S_r \wr S_s \le S_n$.
    There is a bijection between the $(K,S_i\times S_{n-i})$ double cosets of $S_n$ and $P_{r,s}(i)$, the set of partitions of $i$ into at most $s$ parts, all of size at most $r$. 
\end{lemma}

\begin{proof}
    To describe the bijection from $K\backslash S_n/(S_i \times S_{n-i})$ to $P_{r,s}(i)$ explicitly,
    first fix the underlying decomposition $\{1, \ldots, n\} = \{1,\ldots,i\} \cup \{i+1,\ldots, n\}$ for the action of $S_i\times S_{n-i}$. 
    The left cosets in $S_n/(S_i \times S_{n-i})$ are clearly in bijection with the subsets of size $i$ in $\{1, \ldots, n\}$:
    \[
        g(S_i \times S_{n-i}) \longleftrightarrow g(\{1, \ldots, i\}).
    \]
    Now fix a decomposition for the action of $K \cong S_r \wr S_s$: 
    $\{1, \ldots, n\} = B_1 \cup \ldots \cup B_s$,
    where $B_j := \{(j-1)r+1,(j-1)r+2,\ldots,(j-1)r+r\}$ $(j = 1,\ldots,s)$. 
    The elements of $S_s$ permute these blocks,
    and each of the $s$ copies of $S_r$ acts on one of the blocks.
    Given $g \in S_n$, we map the double coset $K g (S_i \times S_{n-i})$ to the partition $\gamma$ which is the non-increasing rearrangement of the sequence
    \[
        (\card{B_1 \cap g(\{1,\ldots,i\})}, \ldots, \card{B_s \cap g(\{1,\ldots,i\})}).
    \]
    This sequence consists of $s$ non-negative integers, each at most $r$, which sum up to $i$. Thus $\gamma \in P_{r,s}(i)$.
    We will show that this map is a bijection.

    For arbitrary $x\in K$ and $y\in S_i\times S_{n-i}$ we have,
    for each $1 \le j \le s$,
    \[
        \card{B_j \cap xgy(\{1,\ldots,i\})} 
        = \card{B_j \cap xg(\{1,\ldots,i\})}
        = \card{x^{-1}(B_j) \cap g(\{1,\ldots,i\})}.
    \]
    The element $x^{-1} \in K = S_r \wr S_s$ permutes the blocks and permutes the elements of each block. This shows that the mapping $Kg(S_i \times S_{n-i}) \mapsto \gamma$ is well defined.

    The mapping from $K\backslash S_n/(S_i \times S_{n-i})$ to $P_{r,s}(i)$ is clearly onto, since
    for each partition $\gamma = (a_1,\ldots,a_s) \in P_{r,s}(i)$ there exists a permutation $g \in S_n$ such that  $\card{B_j \cap g(\{1,\ldots,i\})} = a_j$ $(1 \le j \le s)$. 
		
    Finally, if $K g_1 (S_i \times S_{n-i})$ and $K g_2 (S_i \times S_{n-i})$ are mapped to the same partition~$\gamma$, then there exists a permutation $\pi \in S_s$ satisfying
    \[
        \card{B_j \cap g_1(\{1,\ldots,i\})}
        = \card{B_{\pi(j)} \cap g_2(\{1,\ldots,i\})}
        \qquad (1 \le j \le s).
    \]
    Therefore there exist an element $x \in K = S_r \wr S_s$ such that
    \[
        B_j \cap g_1(\{1,\ldots,i\})
        = B_j \cap xg_2(\{1,\ldots,i\})
        \qquad (1 \le j \le s).
    \]
    It follows that
    \[
        g_1(\{1,\ldots,i\}) = xg_2(\{1,\ldots,i\}),
    \]
    and therefore $g_1 = xg_2y$ for a suitable permutation $y \in S_i \times S_{n-i}$.
\end{proof}

We are now ready to prove Theorem~\ref{prop:e_i_formuli}. 
The proof applies Lemma~\ref{lemma:parametrization} and the explicit bijection described in its proof, combined with Lemma~\ref{lem:char_values_and_inner_product}. 

\begin{proof}[Proof of  Theorem~\ref{prop:e_i_formuli}.]
    Recall that $n=rs$, 
    $K=K_{r,s}\cong S_r\wr S_s\leq S_n$ and 
    $\psi = \phi\uparrow_K^{S_n}$. 
    By Frobenius reciprocity (twice) and 
    Mackey's formula~\cite[(5.2), Problem (5.6)]{Isaacs},
    \begin{equation}\label{Mackey_sum}
    \begin{split}
        e_i 
        &= \langle \psi, \chi^{(1^i) \oplus (n-i)} \rangle 
        = \langle \phi\uparrow_K^{S_n}, (\varepsilon_{S_i} \times 1_{S_{n-i}})\uparrow_{S_i \times S_{n-i}}^{S_n} \rangle \\
        &= \langle {\phi\uparrow_K^{S_n}}\downarrow_{S_i \times S_{n-i}}^{S_n}, \varepsilon_{S_i}\times 1_{S_{n-i}} \rangle \\
        &= 
        \sum_{[g] \in K \backslash S_n / (S_i \times S_{n-i})}
        \langle \phi^g\downarrow^{K^g}_{K^g \cap (S_i \times S_{n-i})}\uparrow_{K^g \cap (S_i \times S_{n-i})}^{S_i \times S_{n-i}}, \varepsilon_{S_i} \times 1_{S_{n-i}} \rangle \\
        &=
        \sum_{[g] \in K \backslash S_n / (S_i \times S_{n-i})}
        \langle \phi^g\downarrow^{K^g}_{K^g \cap (S_i \times S_{n-i})}, (\varepsilon_{S_i} \times 1_{S_{n-i}})\downarrow^{S_i \times S_{n-i}}_{K^g \cap  (S_i \times S_{n-i})} \rangle.
    \end{split}
    \end{equation}
		
    The above sums are indexed by the $(K,S_i\times S_{n-i})$ double cosets in $S_n$. 
    For each representative $g$ of a double coset, 
    $K^g := g^{-1}Kg$ is the corresponding conjugate of $K$ 
    and $\phi^g$ is the character on $K^g$ defined by
    $\phi^g(g^{-1}kg) := \phi(k)$ for all $k \in K$.

    By Lemma~\ref{lemma:parametrization}, 
    these double cosets are parametrized by the partitions in $P_{r,s}(i)$.  Let us determine the summand of \eqref{Mackey_sum} 
    corresponding to a partition $\gamma\in P_{r,s}(i)$ in which part $j$ occurs with multiplicity $k_j = k_j(\gamma)$ $(0 \le j \le r)$. 
    By the bijection described in the proof of Lemma~\ref{lemma:parametrization}, 
    \[
        k_j = \card{\{t \,:\, \card{B_t \cap g(\{1, \ldots, i\})} = j\}}
        \qquad (0 \le j \le r),
    \]
    where $B_t := \{(t-1)r+1,\ldots,(t-1)r+r\}$ $(1 \le t \le s)$. Thus
    \[
        K^g \cap (S_i \times S_{n-i})
        \cong (R_{r,0} \wr S_{k_0}) \times (R_{r,1} \wr S_{k_1}) \times \cdots \times (R_{r,r}\wr S_{k_r}),
    \]
    where $R_{r,j} = S_j \times S_{r-j}$, as above. In particular,
    \begin{align*}
        (\varepsilon_{S_i} \times 1_{S_{n-i}})\downarrow^{S_i \times S_{n-i}}_{K^g \cap (S_i \times S_{n-i})}
        &= \bigotimes_j (\varepsilon_{S_{jk_j}} \times 1_{S_{(r-j)k_j}})\downarrow^{S_{jk_j} \times S_{(r-j)k_j}}_{R_{r,j}\wr S_{k_j}}
        = \bigotimes_j \nu_{r,j,k_j}.
    \end{align*}
    Note that, by Observation~\ref{obs:phi_product}, $\phi_{r,s}$ factors similarly.
    Therefore the corresponding summand in~\eqref{Mackey_sum} is
    \begin{equation}\label{Mackey_summand}
        \langle \phi^g\downarrow^{K^g}_{K^g \cap (S_i \times S_{n-i})},
        (\varepsilon_{S_i} \times 1_{S_{n-i}})\downarrow^{S_i \times S_{n-i}}_{K^g \cap (S_i \times S_{n-i})} \rangle 
        = \prod_{j=0}^{r} \langle \phi_{r,k_j}\downarrow^{S_r\wr S_{k_j}}_{R_{r,j}\wr S_{k_j}}, \nu_{r,j,k_j} \rangle.
    \end{equation}
		
    We have already computed these inner products in Lemma~\ref{lem:char_values_and_inner_product}.
    By \eqref{Mackey_sum}, \eqref{Mackey_summand}, Lemma~\ref{lem:char_values_and_inner_product}
    and Definition~\ref{def:f} we obtain
    \[
        e_{i,(r^s)} 
        = \langle \phi\uparrow_K^{S_n}, (\varepsilon_{S_i} \times 1_{S_{n-i}})\uparrow_{S_i \times S_{n-i}}^{S_n} \rangle
        = \sum_{\gamma\in P_{r,s}(i)} \prod_{j=0}^{r} 
        (-1)^{(j+1)k_j(\gamma)} \binom{(-1)^{j+1} f_j(r)}{k_j(\gamma)},
    \]
    as claimed.
    If $s=1$ then $P_{r,1}(i)$ contains (for $0 \le i \le r$) 
    a unique partition $\gamma = (i)$, for which $k_i(\gamma) = 1$ is the unique nonzero multiplicity. Thus, in this case,  
    \[
        e_{i,(r)} 
        = (-1)^{i+1} \binom{(-1)^{i+1}f_i(r)}{1} 
        = f_i(r).
        \qedhere
    \]
\end{proof}
 
	
\subsection{A product formula}
\label{sec:prod}
	
Now we derive a restatement of Theorem~\ref{prop:e_i_formuli} as a product formula for formal power series. 

\begin{corollary}\label{power_series_form}
    For a positive integer $r$, define the formal power series
    \[
        E_r(x,y) := \sum_{i,s \ge 0} e_{i,(r^s)} x^i y^s.
    \]
    Then
    \[
        E_r(x,y) 
        = \prod_{j=0}^{r} (1 - (-x)^j y)^{(-1)^{j+1}f_j(r)}.
    \]
\end{corollary}
	
\begin{proof}
    Recall the following formal power series expansion, valid for any integer $f$:
    \[
        (1+t)^{f} = \sum_{n=0}^\infty \binom{f}{n} t^n.
    \]
    By Theorem~\ref{prop:e_i_formuli},
    with the obvious extension for $s=0$,
    \begin{align*}
        E_r(x,y) 
        &= \sum_{i,s \ge 0} \left(
        \sum_{\substack{k_0,\ldots,k_r \ge 0 \\ \sum_j k_j=s \\ \sum_j jk_j = i}} 
        \prod_{j = 0}^{r}
        (-1)^{(j+1)k_j} \binom{(-1)^{j+1} f_j(r)}{k_j}
        \right) x^i y^s \\
        &= \prod_{j=0}^{r} \sum_{k_j=0}^\infty
        (-1)^{(j+1)k_j} \binom{(-1)^{j+1} f_j(r)}{k_j}
        x^{jk_j} y^{k_j} \\
        &= \prod_{j=0}^{r} (1 + (-1)^{j+1} x^j y)^{(-1)^{j+1}f_j(r)},
    \end{align*}
    as required.
\end{proof}
	
\begin{remark}\label{rem:le_3}
    For small $r$ and arbitrary $s$, Corollary~\ref{power_series_form}  enables us to determine explicitly the hook-multiplicities $m_{i,(r^s)}$. 
    This is done by recalling Equation~\eqref{eq:alternating} and the fact that, by definition, $e_{i,(r^s)}$ is the coefficient of $x^i y^s$ in $E_r(x,y)$. 
    For example, by Corollary~\ref{power_series_form} and Observation~\ref{obs:f_values}, 
    $E_2(x,y) = E_3(x,y) = (1+xy)(1-x^2y)^{-1}$. 
    Thus, for $r \in \{2,3\}$ and any $s \ge 1$, the value of $e_{i,(r^s)}$ is $1$ for $i \in \{2s-1,2s\}$ and zero otherwise. 
    Combining this with Equation~\eqref{eq:alternating}, 
    it follows that, for $r \in \{2,3\}$ and $s \ge 1$, the hook multiplicity $m_{i,(r^s)}$ 
    is 1 for $i = 2s-1$ and zero otherwise. 
\end{remark}	


\subsection{Non-negativity}\label{sec:nn}
		
Now we are ready to prove Theorem~\ref{t:main_hook-alternating2}.
	
\begin{proof}[Proof of Theorem~\ref{t:main_hook-alternating2}.]
    Assume that $r$ is not square-free.
	Recall, from Definition~\ref{def:F_polynomial}, the notation $F_r(x) := \sum_{j=0}^{r} f_j(r) x^j$.
	By 	Proposition~\ref{prop:cycles_nonnegativity} and Corollary~\ref{t:F_and_M}	
    we may write $F_r(x) = (1+x)^2 G_r(x)$, where $G_r(x) = \sum_{j=0}^{r-2} g_j(r) x^j$ is a polynomial with non-negative integer coefficients. 
	Let $g_j(r) := 0$ for $j < 0$ or $j > r-2$. Then
	\[
	   f_j(r) = g_j(r) + 2g_{j-1}(r) + g_{j-2}(r)
	   \qquad (\forall j).
	\]
    Therefore, by Corollary~\ref{power_series_form},
    \begin{equation*}\begin{split}
        E_r(x,y) 
        &= \prod_{j \ge 0} (1-(-x)^j y)^{(-1)^{j+1}f_j(r)} \\
        &= \prod_{j \ge 0} (1-(-x)^j y)^{(-1)^{j+1}(g_j(r) + 2g_{j-1}(r) + g_{j-2}(r))} \\
        &= \prod_{j \ge 0} \left(
        \frac{(1 - (-x)^{j} y)(1 - (-x)^{j+2} y)}{(1 - (-x)^{j+1} y)^2} \right)^{(-1)^{j+1} g_j(r)}.
    \end{split}
    \end{equation*}
    We claim that each factor in 
    this product has the form $1 + (1+x)^2 p_j(x,y)$, with $p_j(x,y)$ a formal power series with non-negative integer coefficients. 
    This implies that $(E_r(x,y)-1)/(1+x)^2$ is itself a formal power series with non-negative integer coefficients, completing the proof of Theorem~\ref{t:main_hook-alternating2}.
	
    Indeed, if $j$ is odd then the corresponding factor is
    \[
        \left( \frac{(1 + x^{j} y)(1 + x^{j+2} y)}{(1 - x^{j+1} y)^2} \right)^{g_j(r)}
        = \left( 1 + \frac{(x+1)^2 x^j y}{(1 - x^{j+1} y)^2} \right)^{g_j(r)},
    \]
    where
    \[
        \frac{x^j y}{(1 - x^{j+1} y)^2}
        = x^j y \cdot \sum_{i \ge 0} (i+1) (x^{j+1} y)^i
    \]
    is a formal power series with non-negative integer coefficients.
	
    Finally, if $j$ is even then the corresponding factor is
    \[
        \left( \frac{(1 + x^{j+1} y)^2}{(1 - x^{j} y)(1 - x^{j+2} y)} \right)^{g_j(r)}
        = \left( 1 + \frac{(x+1)^2 x^j y}{(1 - x^{j} y)(1 - x^{j+2} y)} \right)^{g_j(r)},
    \]
    where
    \[
        \frac{x^j y}{(1 - x^{j} y)(1 - x^{j+2} y)}
        = x^j y \cdot \sum_{i \ge 0} (x^{j} y)^i \cdot \sum_{k \ge 0} (x^{j+2} y)^k
    \]
    is a formal power series with non-negative integer coefficients.
\end{proof}
